% ----------------------------------------------------------
\section{Considerações Finais}\label{sec:consideracoesFinais}

Este trabalho comparou duas estratégias para o Problema de Coloração de Grafos com três cores: uma solução original por força bruta e uma versão otimizada baseada em redução para SAT, conversão para 3-CNF-SAT e DPLL com propagação unitária. A análise confirma que a enumeração direta tem valor como linha de base, pois sua correção decorre da cobertura completa do espaço de colorações e sua memória adicional é reduzida. O custo exponencial, entretanto, aparece de forma imediata quando a instância não admite coloração, já que todas as $3^{|V|}$ candidatas precisam ser rejeitadas.

A versão otimizada foi efetiva nas instâncias avaliadas porque a codificação booleana tornou explícitas as restrições do problema. As cláusulas de unicidade de cor e de incompatibilidade entre vértices adjacentes permitem que o DPLL derive consequências antes de completar uma coloração. Em grafos completos não 3-coloríveis, essa representação produziu contradições com poucas decisões, reduzindo o tempo de execução em relação à força bruta. A melhora, porém, não equivale a uma mudança de classe de complexidade: 3-coloração e SAT permanecem problemas de decisão exponenciais no pior caso, e \textcite{Silva2026PAA13} situam 3-CNF-SAT entre as reduções centrais para NP-Completude.

O principal custo da otimização foi o aumento de memória. A solução SAT/DPLL precisa armazenar variáveis, cláusulas, mapeamentos e atribuições parciais; por isso, a média de memória observada superou a da força bruta. Esse trade-off é aceitável quando o tempo de busca domina o custo total, especialmente em grafos densos ou em instâncias insatisfazíveis. Em grafos muito pequenos, esparsos ou com coloração encontrada nas primeiras tentativas, o custo fixo da redução pode neutralizar o ganho obtido pela propagação. A implementação também usa uma conversão 3-CNF adequada ao caso $k=3$, mas limitada para valores maiores de $k$, nos quais cláusulas longas exigiriam transformação geral com variáveis auxiliares.

As ameaças à validade concentram-se no tamanho e na diversidade das instâncias. Os experimentos usaram grafos completos entre três e nove vértices e uma repetição por tamanho, o que favorece a observação de conflitos fortes e não representa todos os padrões estruturais encontrados em aplicações reais. Além disso, as medições foram feitas em Python, com \texttt{tracemalloc} para memória e tempos muito pequenos, sensíveis a variações do ambiente de execução. Estudos posteriores devem incluir grafos aleatórios com diferentes densidades, grafos planares, instâncias satisfazíveis e insatisfazíveis balanceadas, múltiplas repetições e comparação com heurísticas clássicas de coloração, como ordenação por grau e retrocesso com poda.

Como recomendação prática, a solução otimizada é mais vantajosa quando há densidade de restrições suficiente para alimentar a propagação unitária ou quando a inexistência de coloração precisa ser demonstrada sem enumerar colorações completas. A força bruta permanece útil para validação, ensino e testes pequenos, nos quais simplicidade e baixo uso de memória são preferíveis à infraestrutura de redução. A comparação mostra que reduções polinomiais para SAT são ferramentas relevantes não apenas para provas de complexidade, mas também para engenharia de algoritmos: elas deslocam o problema para uma representação na qual inferências locais podem reduzir substancialmente a busca, desde que o custo de memória e as limitações da codificação sejam controlados.
